\section{量子化族と多項式不変量}
\label{sec:primitive-setup-v36}

基礎体を $\mathbb C$ とする。整数 $j,r\ge0$、$m=2j+r\ge3$ に対して
\[
 R_m=\mathbb C[x,y,z,w],\quad \deg(x,y,z,w)=(1,1,1,m),\quad
 q=xz-y^2,\quad d=m+r+1
\]
と置く。次の二つの導分を使う。
\begin{align*}
 N&=x\partial_y+2y\partial_z-q^jz^r\partial_w,\\
 X_\alpha&=\alpha x\partial_x+(\alpha+1)y\partial_y+(\alpha+2)z\partial_z
                  +(m\alpha+d)w\partial_w.
\end{align*}
$N$ は局所冪零、$[N,X_\alpha]=N$ である。
同じベクトル空間に有限積
\[
 a\star b=\sum_{n\ge0}s^nN^n(a)\binom{X_\alpha}{n}(b)
\]
を入れた代数を $\mathcal S_{j,r}(\alpha,s)$ または
$\mathcal S^{(q)}_{j,r}(\alpha,s)$ と書く。
積の結合性はCGG恒等式による\cite[Proposition~2.2]{LecoutreSierra2019}。
本稿では量子化係数 $s\ne0$ を固定する。

\begin{proposition}[基本的な入力]
\label{prop:primitive-ring-input-v36}
$\mathcal S$ はnoetherianなAS正則四次元整域であり、
\[
 H_{\mathcal S}(t)=\frac1{(1-t)^3(1-t^m)}.
\]
$T=1+sN$ とその全スカラー冪は自己同型であり、Nakayama自己同型は
$T^{(m+3)\alpha+m+r+3}$ である。
従って通常のCalabi--Yau代表は
$\alpha=-(m+r+3)/(m+3)$ で与えられる。
\end{proposition}

これらの性質と以下の不変環表示は、公表稿IX\cite{VKPoisson34}の反復Ore表示と
多項式不変量の計算による。量子化の発散だけからCalabi--Yau性を推論しない。
\[
 c_r=(-1)^r\frac{(2r)!!}{(2r+1)!!},\qquad
 J=x^{r+1}w+q^j\sum_{k=0}^r\binom rk\frac{q^{r-k}y^{2k+1}}{2k+1},\qquad
 \Delta=\frac{J^2+c_r^2q^d}{x^{r+1}}
\]
と置くと、$\Delta$ は多項式であり、
\[
 \ker N=\mathbb C[x,q,J,\Delta]
 \simeq\mathbb C[X,Q,U,V]/(X^{r+1}V-U^2-c_r^2Q^d).
\]
$N$ 不変元同士の星積は通常の可換積と等しい。
通常CY値での全中心は
\[
 Z(\mathcal S)=
 \begin{cases}
 \mathbb C&r\ge2,\\
 \mathbb C[q,J]&r=0,\\
 \mathbb C[\mathsf S,\mathsf P,\mathsf R]/
 (\mathsf R^2-\mathsf S\mathsf P^2-\frac{16}{81}\mathsf S^{m+3})&r=1,
 \end{cases}
\]
ただし最後の場合は
\[
 \mathsf S=q^2\Delta,\quad
 \mathsf P=2xJ\Delta^{(m+3)/2},\quad
 \mathsf R=q(x^2\Delta+J^2)\Delta^{(m+3)/2}
\]
である。この中心の全体性には、正規元 $x$ による飽和と零重量単項式の分類を使う。

以後、$\mathcal S[x^{-1}]$ の中心、射影的次数零chartの中心、
および $Z(\mathcal S)$ は区別して記す。
また、primitive idealとは単純左加群の零化イデアルを指し、
$\operatorname{Prim}\mathcal S$ には両側イデアルの包含で定めるJacobson位相を入れる。
本稿の全primitive idealは完全素となるので、各商環の分数斜体を通常のOre斜体として扱える。
